The rapid expansion of the Internet of Things (IoT) in industrial settings has created a critical need for automated anomaly detection. This report presents the design and implementation of a robust IoT Anomaly Detection System that utilizes unsupervised machine learning to identify irregularities in sensor data. The system features a comprehensive 9-stage data cleaning pipeline that handles noise, missing values, and statistical outliers using the Interquartile Range (IQR) method. Three machine learning algorithms—Isolation Forest, One-Class SVM, and K-Means—are implemented and compared to provide a versatile detection framework. Furthermore, the system incorporates diagnostic metadata, such as network delay, to provide a holistic view of both data integrity and infrastructure health. A professional web-based dashboard was developed using Angular and .NET to allow for interactive data exploration, multi-model comparison, and professional-grade reporting. The results demonstrate that combining advanced preprocessing with unsupervised models significantly improves the reliability of anomaly detection in industrial IoT environments. However, the system may face challenges in detecting complex non-linear anomalies in high-noise environments, a limitation that can be overcome by incorporating hybrid ensemble learning techniques. The evaluation results show that the Isolation Forest model achieved the highest accuracy of 92.81\% , while K-Means and One-Class SVM followed with accuracies of 92.65\% and 92.44\%, respectively. The superior performance of the Isolation Forest model is attributed to its robustness in isolating anomalies within the complex, high-variance patterns typical of industrial sensor data.

\vspace{8pt}
\textbf{Keywords:} IoT, Anomaly Detection, Machine Learning, Unsupervised Learning, Isolation Forest, Data Cleaning, Industrial Sensors.
